% =====================================================================
%  Anatomía de un LLM — Fase I — Módulo 3
%  Álgebra lineal II: rango, autovalores, SVD y PCA
%  Compilar: latexmk -r ../latexmkrc M03-algebra-lineal-2.tex
% =====================================================================
\documentclass[11pt,a4paper,twoside,openany]{book}
\usepackage{../estilo-apuntes}

\newcommand{\ApuntesCurso}{Anatomía de un LLM}
\newcommand{\ApuntesModulo}{Módulo 3}

% --- Notación propia del módulo --------------------------------------
\DeclareMathOperator{\rango}{rango}
\DeclareMathOperator{\traza}{tr}
\DeclareMathOperator{\gen}{gen}
\DeclareMathOperator{\diag}{diag}
\DeclareMathOperator{\Var}{Var}
\newcommand{\Col}{\mathcal{C}}
\newcommand{\Nuc}{\mathcal{N}}
\newcommand{\T}{^{\top}}
\newcommand{\dmodel}{d_{\mathrm{model}}}
\usetikzlibrary{calc}
\usepackage{xurl}

\hypersetup{pdftitle={Módulo 3. Álgebra lineal II (resumen)},pdfauthor={Anatomía de un LLM}}

\makeatletter
\newlabel{ej:coordenadas}{{1.9}{2}{Cambio de coordenadas}{ejemplo.9}{}}
\newlabel{ej:coordenadas@cref}{{[ejemplo][9][]1.9}{[1][2][]2}}
\newlabel{ej:M}{{1.15}{3}{}{ejemplo.15}{}}
\newlabel{ej:M@cref}{{[ejemplo][15][]1.15}{[1][3][]3}}
\newlabel{eje:1-vf}{{1.1}{4}{}{ejercicio.1.1}{}}
\newlabel{eje:1-vf@cref}{{[ejercicio][1][1]1.1}{[1][4][]4}}
\newlabel{eje:1-rango1}{{1.2}{4}{}{ejercicio.1.2}{}}
\newlabel{eje:1-rango1@cref}{{[ejercicio][2][1]1.2}{[1][4][]4}}
\newlabel{eje:1-mml}{{1.3}{4}{}{ejercicio.1.3}{}}
\newlabel{eje:1-mml@cref}{{[ejercicio][3][1]1.3}{[1][4][]4}}
\newlabel{eje:1-u4}{{1.4}{4}{}{ejercicio.1.4}{}}
\newlabel{eje:1-u4@cref}{{[ejercicio][4][1]1.4}{[1][4][]4}}
\newlabel{eje:1-nucleo}{{1.5}{4}{}{ejercicio.1.5}{}}
\newlabel{eje:1-nucleo@cref}{{[ejercicio][5][1]1.5}{[1][4][]4}}
\newlabel{eje:1-producto}{{1.6}{4}{}{ejercicio.1.6}{}}
\newlabel{eje:1-producto@cref}{{[ejercicio][6][1]1.6}{[1][4][]4}}
\newlabel{ej:inversa}{{2.5}{6}{}{ejemplo.5}{}}
\newlabel{ej:inversa@cref}{{[ejemplo][5][]2.5}{[1][6][]6}}
\newlabel{ej:giro}{{2.13}{7}{Giros del plano}{ejemplo.13}{}}
\newlabel{ej:giro@cref}{{[ejemplo][13][]2.13}{[1][7][]7}}
\newlabel{eje:2-equiv}{{2.1}{8}{}{ejercicio.2.1}{}}
\newlabel{eje:2-equiv@cref}{{[ejercicio][1][2]2.1}{[1][8][]8}}
\newlabel{eje:2-inversa}{{2.2}{8}{}{ejercicio.2.2}{}}
\newlabel{eje:2-inversa@cref}{{[ejercicio][2][2]2.2}{[1][8][]8}}
\newlabel{eje:2-mml}{{2.3}{8}{}{ejercicio.2.3}{}}
\newlabel{eje:2-mml@cref}{{[ejercicio][3][2]2.3}{[1][8][]8}}
\newlabel{eje:2-giros}{{2.4}{8}{}{ejercicio.2.4}{}}
\newlabel{eje:2-giros@cref}{{[ejercicio][4][2]2.4}{[1][8][]8}}
\newlabel{eje:2-haar}{{2.5}{8}{}{ejercicio.2.5}{}}
\newlabel{eje:2-haar@cref}{{[ejercicio][5][2]2.5}{[1][8][]8}}
\newlabel{eje:2-isometria}{{2.6}{9}{}{ejercicio.2.6}{}}
\newlabel{eje:2-isometria@cref}{{[ejercicio][6][2]2.6}{[1][9][]9}}
\newlabel{ej:S}{{3.3}{10}{}{ejemplo.3}{}}
\newlabel{ej:S@cref}{{[ejemplo][3][]3.3}{[1][10][]10}}
\newlabel{ej:potencia}{{3.12}{13}{}{ejemplo.12}{}}
\newlabel{ej:potencia@cref}{{[ejemplo][12][]3.12}{[1][13][]13}}
\newlabel{eje:3-mml}{{3.1}{14}{}{ejercicio.3.1}{}}
\newlabel{eje:3-mml@cref}{{[ejercicio][1][3]3.1}{[1][13][]14}}
\newlabel{eje:3-autovec}{{3.2}{14}{}{ejercicio.3.2}{}}
\newlabel{eje:3-autovec@cref}{{[ejercicio][2][3]3.2}{[1][14][]14}}
\newlabel{eje:3-u2}{{3.3}{14}{}{ejercicio.3.3}{}}
\newlabel{eje:3-u2@cref}{{[ejercicio][3][3]3.3}{[1][14][]14}}
\newlabel{eje:3-u3}{{3.4}{14}{}{ejercicio.3.4}{}}
\newlabel{eje:3-u3@cref}{{[ejercicio][4][3]3.4}{[1][14][]14}}
\newlabel{eje:3-reales}{{3.5}{14}{}{ejercicio.3.5}{}}
\newlabel{eje:3-reales@cref}{{[ejercicio][5][3]3.5}{[1][14][]14}}
\newlabel{eje:3-potencia}{{3.6}{14}{}{ejercicio.3.6}{}}
\newlabel{eje:3-potencia@cref}{{[ejercicio][6][3]3.6}{[1][14][]14}}
\newlabel{fig:svd}{{4.1}{16}{Interpretación geométrica de $A=U\Sigma V\T $ para $m=n=2$. La imagen de la circunferencia unidad es una elipse cuyos semiejes miden $\sigma _1$ y $\sigma _2$ y apuntan en las direcciones $u_1$ y $u_2$}{figure.4.1}{}}
\newlabel{fig:svd@cref}{{[figure][1][4]4.1}{[1][15][]16}}
\newlabel{ej:svd23}{{4.4}{16}{}{ejemplo.4}{}}
\newlabel{ej:svd23@cref}{{[ejemplo][4][]4.4}{[1][16][]16}}
\newlabel{eje:4-vf}{{4.1}{17}{}{ejercicio.4.1}{}}
\newlabel{eje:4-vf@cref}{{[ejercicio][1][4]4.1}{[1][17][]17}}
\newlabel{eje:4-u1}{{4.2}{17}{}{ejercicio.4.2}{}}
\newlabel{eje:4-u1@cref}{{[ejercicio][2][4]4.2}{[1][17][]17}}
\newlabel{eje:4-u5}{{4.3}{18}{}{ejercicio.4.3}{}}
\newlabel{eje:4-u5@cref}{{[ejercicio][3][4]4.3}{[1][18][]18}}
\newlabel{eje:4-mml}{{4.4}{18}{}{ejercicio.4.4}{}}
\newlabel{eje:4-mml@cref}{{[ejercicio][4][4]4.4}{[1][18][]18}}
\newlabel{eje:4-frobenius}{{4.5}{18}{}{ejercicio.4.5}{}}
\newlabel{eje:4-frobenius@cref}{{[ejercicio][5][4]4.5}{[1][18][]18}}
\newlabel{eje:4-autovalores}{{4.6}{18}{}{ejercicio.4.6}{}}
\newlabel{eje:4-autovalores@cref}{{[ejercicio][6][4]4.6}{[1][18][]18}}
\newlabel{eje:4-espectral}{{4.7}{18}{}{ejercicio.4.7}{}}
\newlabel{eje:4-espectral@cref}{{[ejercicio][7][4]4.7}{[1][18][]18}}
\newlabel{ej:A1}{{5.3}{20}{}{ejemplo.3}{}}
\newlabel{ej:A1@cref}{{[ejemplo][3][]5.3}{[1][20][]20}}
\newlabel{fig:lora}{{5.1}{21}{Capa lineal con LoRA. La entrada atraviesa la matriz congelada $W_0$ y, en paralelo, el cuello de dimensión $r$ formado por $A$ y $B$; la corrección $\Delta W=BA$ tiene rango a lo sumo~$r$}{figure.5.1}{}}
\newlabel{fig:lora@cref}{{[figure][1][5]5.1}{[1][20][]21}}
\newlabel{ej:pca}{{5.7}{22}{}{ejemplo.7}{}}
\newlabel{ej:pca@cref}{{[ejemplo][7][]5.7}{[1][22][]22}}
\newlabel{eje:5-centrar}{{5.1}{23}{}{ejercicio.5.1}{}}
\newlabel{eje:5-centrar@cref}{{[ejercicio][1][5]5.1}{[1][22][]23}}
\newlabel{eje:5-varexp}{{5.2}{23}{}{ejercicio.5.2}{}}
\newlabel{eje:5-varexp@cref}{{[ejercicio][2][5]5.2}{[1][23][]23}}
\newlabel{eje:5-svd}{{5.3}{23}{}{ejercicio.5.3}{}}
\newlabel{eje:5-svd@cref}{{[ejercicio][3][5]5.3}{[1][23][]23}}
\newlabel{eje:5-mml}{{5.4}{23}{}{ejercicio.5.4}{}}
\newlabel{eje:5-mml@cref}{{[ejercicio][4][5]5.4}{[1][23][]23}}
\newlabel{eje:5-pca}{{5.5}{23}{}{ejercicio.5.5}{}}
\newlabel{eje:5-pca@cref}{{[ejercicio][5][5]5.5}{[1][23][]23}}
\newlabel{eje:5-coste}{{5.6}{23}{}{ejercicio.5.6}{}}
\newlabel{eje:5-coste@cref}{{[ejercicio][6][5]5.6}{[1][23][]23}}
\newlabel{eje:5-lora}{{5.7}{23}{}{ejercicio.5.7}{}}
\newlabel{eje:5-lora@cref}{{[ejercicio][7][5]5.7}{[1][23][]23}}
\newlabel{fig:formas}{{6.1}{26}{Transformaciones de forma en un LLM: búsqueda en la tabla de embeddings, capa lineal, puntuaciones de atención y logits con pesos atados}{figure.6.1}{}}
\newlabel{fig:formas@cref}{{[figure][1][6]6.1}{[1][25][]26}}
\newlabel{eje:6-formas}{{6.1}{27}{}{ejercicio.6.1}{}}
\newlabel{eje:6-formas@cref}{{[ejercicio][1][6]6.1}{[1][27][]27}}
\newlabel{eje:6-d2l}{{6.2}{27}{}{ejercicio.6.2}{}}
\newlabel{eje:6-d2l@cref}{{[ejercicio][2][6]6.2}{[1][27][]27}}
\newlabel{eje:6-embeddings}{{6.3}{27}{}{ejercicio.6.3}{}}
\newlabel{eje:6-embeddings@cref}{{[ejercicio][3][6]6.3}{[1][27][]27}}
\newlabel{eje:6-atencion}{{6.4}{27}{}{ejercicio.6.4}{}}
\newlabel{eje:6-atencion@cref}{{[ejercicio][4][6]6.4}{[1][27][]27}}
\newlabel{eje:6-orden}{{6.5}{27}{}{ejercicio.6.5}{}}
\newlabel{eje:6-orden@cref}{{[ejercicio][5][6]6.5}{[1][27][]27}}
\newlabel{eje:6-cuello}{{6.6}{27}{}{ejercicio.6.6}{}}
\newlabel{eje:6-cuello@cref}{{[ejercicio][6][6]6.6}{[1][27][]27}}
\newlabel{ap:soluciones}{{A}{28}{Soluciones de los ejercicios}{appendix.A}{}}
\newlabel{ap:soluciones@cref}{{[appendix][1][2147483647]A}{[1][28][]28}}
\newlabel{sol:eje:1-vf}{{A.1}{28}{Tema 1}{section*.3}{}}
\newlabel{sol:eje:1-vf@cref}{{[subappendix][1][2147483647,1]A.1}{[1][28][]28}}
\newlabel{sol:eje:1-rango1}{{A.1}{28}{Tema 1}{section*.4}{}}
\newlabel{sol:eje:1-rango1@cref}{{[subappendix][1][2147483647,1]A.1}{[1][28][]28}}
\newlabel{sol:eje:1-mml}{{A.1}{29}{Tema 1}{section*.5}{}}
\newlabel{sol:eje:1-mml@cref}{{[subappendix][1][2147483647,1]A.1}{[1][29][]29}}
\newlabel{sol:eje:1-u4}{{A.1}{29}{Tema 1}{section*.6}{}}
\newlabel{sol:eje:1-u4@cref}{{[subappendix][1][2147483647,1]A.1}{[1][29][]29}}
\newlabel{sol:eje:1-nucleo}{{A.1}{30}{Tema 1}{section*.7}{}}
\newlabel{sol:eje:1-nucleo@cref}{{[subappendix][1][2147483647,1]A.1}{[1][30][]30}}
\newlabel{sol:eje:1-producto}{{A.1}{30}{Tema 1}{section*.8}{}}
\newlabel{sol:eje:1-producto@cref}{{[subappendix][1][2147483647,1]A.1}{[1][30][]30}}
\newlabel{sol:eje:2-equiv}{{A.2}{31}{Tema 2}{section*.9}{}}
\newlabel{sol:eje:2-equiv@cref}{{[subappendix][2][2147483647,1]A.2}{[1][31][]31}}
\newlabel{sol:eje:2-inversa}{{A.2}{31}{Tema 2}{section*.10}{}}
\newlabel{sol:eje:2-inversa@cref}{{[subappendix][2][2147483647,1]A.2}{[1][31][]31}}
\newlabel{sol:eje:2-mml}{{A.2}{32}{Tema 2}{section*.11}{}}
\newlabel{sol:eje:2-mml@cref}{{[subappendix][2][2147483647,1]A.2}{[1][31][]32}}
\newlabel{sol:eje:2-giros}{{A.2}{32}{Tema 2}{section*.12}{}}
\newlabel{sol:eje:2-giros@cref}{{[subappendix][2][2147483647,1]A.2}{[1][32][]32}}
\newlabel{sol:eje:2-haar}{{A.2}{32}{Tema 2}{section*.13}{}}
\newlabel{sol:eje:2-haar@cref}{{[subappendix][2][2147483647,1]A.2}{[1][32][]32}}
\newlabel{sol:eje:2-isometria}{{A.2}{33}{Tema 2}{section*.14}{}}
\newlabel{sol:eje:2-isometria@cref}{{[subappendix][2][2147483647,1]A.2}{[1][33][]33}}
\newlabel{sol:eje:3-mml}{{A.3}{33}{Tema 3}{section*.15}{}}
\newlabel{sol:eje:3-mml@cref}{{[subappendix][3][2147483647,1]A.3}{[1][33][]33}}
\newlabel{sol:eje:3-autovec}{{A.3}{34}{Tema 3}{section*.16}{}}
\newlabel{sol:eje:3-autovec@cref}{{[subappendix][3][2147483647,1]A.3}{[1][34][]34}}
\newlabel{sol:eje:3-u2}{{A.3}{34}{Tema 3}{section*.17}{}}
\newlabel{sol:eje:3-u2@cref}{{[subappendix][3][2147483647,1]A.3}{[1][34][]34}}
\newlabel{sol:eje:3-u3}{{A.3}{35}{Tema 3}{section*.18}{}}
\newlabel{sol:eje:3-u3@cref}{{[subappendix][3][2147483647,1]A.3}{[1][35][]35}}
\newlabel{sol:eje:3-reales}{{A.3}{35}{Tema 3}{section*.19}{}}
\newlabel{sol:eje:3-reales@cref}{{[subappendix][3][2147483647,1]A.3}{[1][35][]35}}
\newlabel{sol:eje:3-potencia}{{A.3}{36}{Tema 3}{section*.20}{}}
\newlabel{sol:eje:3-potencia@cref}{{[subappendix][3][2147483647,1]A.3}{[1][35][]36}}
\newlabel{sol:eje:4-vf}{{A.4}{36}{Tema 4}{section*.21}{}}
\newlabel{sol:eje:4-vf@cref}{{[subappendix][4][2147483647,1]A.4}{[1][36][]36}}
\newlabel{sol:eje:4-u1}{{A.4}{36}{Tema 4}{section*.22}{}}
\newlabel{sol:eje:4-u1@cref}{{[subappendix][4][2147483647,1]A.4}{[1][36][]36}}
\newlabel{sol:eje:4-u5}{{A.4}{37}{Tema 4}{section*.23}{}}
\newlabel{sol:eje:4-u5@cref}{{[subappendix][4][2147483647,1]A.4}{[1][37][]37}}
\newlabel{sol:eje:4-mml}{{A.4}{37}{Tema 4}{section*.24}{}}
\newlabel{sol:eje:4-mml@cref}{{[subappendix][4][2147483647,1]A.4}{[1][37][]37}}
\newlabel{sol:eje:4-frobenius}{{A.4}{38}{Tema 4}{section*.25}{}}
\newlabel{sol:eje:4-frobenius@cref}{{[subappendix][4][2147483647,1]A.4}{[1][38][]38}}
\newlabel{sol:eje:4-autovalores}{{A.4}{38}{Tema 4}{section*.26}{}}
\newlabel{sol:eje:4-autovalores@cref}{{[subappendix][4][2147483647,1]A.4}{[1][38][]38}}
\newlabel{sol:eje:4-espectral}{{A.4}{39}{Tema 4}{section*.27}{}}
\newlabel{sol:eje:4-espectral@cref}{{[subappendix][4][2147483647,1]A.4}{[1][38][]39}}
\newlabel{sol:eje:5-centrar}{{A.5}{39}{Tema 5}{section*.28}{}}
\newlabel{sol:eje:5-centrar@cref}{{[subappendix][5][2147483647,1]A.5}{[1][39][]39}}
\newlabel{sol:eje:5-varexp}{{A.5}{39}{Tema 5}{section*.29}{}}
\newlabel{sol:eje:5-varexp@cref}{{[subappendix][5][2147483647,1]A.5}{[1][39][]39}}
\newlabel{sol:eje:5-svd}{{A.5}{40}{Tema 5}{section*.30}{}}
\newlabel{sol:eje:5-svd@cref}{{[subappendix][5][2147483647,1]A.5}{[1][40][]40}}
\newlabel{sol:eje:5-mml}{{A.5}{40}{Tema 5}{section*.31}{}}
\newlabel{sol:eje:5-mml@cref}{{[subappendix][5][2147483647,1]A.5}{[1][40][]40}}
\newlabel{sol:eje:5-pca}{{A.5}{41}{Tema 5}{section*.32}{}}
\newlabel{sol:eje:5-pca@cref}{{[subappendix][5][2147483647,1]A.5}{[1][41][]41}}
\newlabel{sol:eje:5-coste}{{A.5}{41}{Tema 5}{section*.33}{}}
\newlabel{sol:eje:5-coste@cref}{{[subappendix][5][2147483647,1]A.5}{[1][41][]41}}
\newlabel{sol:eje:5-lora}{{A.5}{42}{Tema 5}{section*.34}{}}
\newlabel{sol:eje:5-lora@cref}{{[subappendix][5][2147483647,1]A.5}{[1][42][]42}}
\newlabel{sol:eje:6-formas}{{A.6}{42}{Tema 6}{section*.35}{}}
\newlabel{sol:eje:6-formas@cref}{{[subappendix][6][2147483647,1]A.6}{[1][42][]42}}
\newlabel{sol:eje:6-d2l}{{A.6}{43}{Tema 6}{section*.36}{}}
\newlabel{sol:eje:6-d2l@cref}{{[subappendix][6][2147483647,1]A.6}{[1][43][]43}}
\newlabel{sol:eje:6-embeddings}{{A.6}{43}{Tema 6}{section*.37}{}}
\newlabel{sol:eje:6-embeddings@cref}{{[subappendix][6][2147483647,1]A.6}{[1][43][]43}}
\newlabel{sol:eje:6-atencion}{{A.6}{44}{Tema 6}{section*.38}{}}
\newlabel{sol:eje:6-atencion@cref}{{[subappendix][6][2147483647,1]A.6}{[1][43][]44}}
\newlabel{sol:eje:6-orden}{{A.6}{44}{Tema 6}{section*.39}{}}
\newlabel{sol:eje:6-orden@cref}{{[subappendix][6][2147483647,1]A.6}{[1][44][]44}}
\newlabel{sol:eje:6-cuello}{{A.6}{44}{Tema 6}{section*.40}{}}
\newlabel{sol:eje:6-cuello@cref}{{[subappendix][6][2147483647,1]A.6}{[1][44][]44}}
\makeatother

\begin{document}
\frontmatter
\begin{titlepage}
\centering
\vspace*{2.5cm}
{\large\scshape Anatomía de un LLM}\par
\vspace{0.4cm}
{\normalsize Fundamentos matemáticos de los modelos de lenguaje}\par
\vspace{2.2cm}
\rule{\textwidth}{0.8pt}\par\vspace{0.5cm}
{\small\scshape Fase I \textperiodcentered\ Fundamentos matemáticos}\par\vspace{0.3cm}
{\Huge\bfseries Módulo 3\par}\vspace{0.3cm}
{\LARGE Álgebra lineal II\par}\vspace{0.2cm}
{\large Rango, autovalores, SVD y PCA\par}
\vspace{0.5cm}\rule{\textwidth}{0.4pt}\par
\vspace{1.6cm}
{\large Resumen del módulo}\par
\vfill
{\small Edición de octubre de 2026}\par
\end{titlepage}

\chapter*{Sobre este resumen}
\addcontentsline{toc}{chapter}{Sobre este resumen}
\markboth{Sobre este resumen}{}

Este documento reúne lo indispensable de los apuntes del Módulo 3: el texto expositivo, todas las definiciones y todos los resultados (teoremas, proposiciones, lemas y corolarios), las tablas y las observaciones de contenido. Se han suprimido las demostraciones, los ejemplos, las lecturas, las figuras, los ejercicios y sus soluciones. La numeración de temas, secciones, enunciados y ecuaciones es la misma que en los apuntes, de modo que cualquier referencia puede buscarse allí; las que remiten a elementos suprimidos (un ejemplo, una figura, una ecuación de una demostración) conservan su número en los apuntes.

\section*{Notación}
\begin{center}
\small
\begin{tabular}{@{}ll@{}}
\toprule
Símbolo & Significado \\
\midrule
$A\in\R^{m\times n}$, $A\T$ & matriz de $m$ filas y $n$ columnas y su traspuesta \\
$\gen\{v_1,\dots,v_k\}$ & subespacio generado por $v_1,\dots,v_k$ \\
$\Col(A)$, $\Nuc(A)$ & espacio columna (imagen) y núcleo de $A$ \\
$\rango(A)$, $\traza(A)$, $\det(A)$ & rango, traza y determinante \\
$I_n$, $e_i$ & matriz identidad y $i$-ésimo vector de la base canónica \\
$\lambda_i$, $\Lambda=\diag(\lambda_1,\dots,\lambda_n)$ & autovalores y matriz diagonal de autovalores \\
$Q$ & matriz ortogonal, $Q\T Q=I$ \\
$A=U\Sigma V\T$, $\sigma_1\ge\sigma_2\ge\dots\ge 0$ & descomposición en valores singulares y valores singulares \\
$A_k=\sum_{i\le k}\sigma_i u_iv_i\T$ & truncamiento de rango $k$ de la SVD \\
$\norm{A}_F$, $\norm{A}_2$ & norma de Frobenius y norma espectral \\
$X\in\R^{n\times d}$, $C=X\T X/n$ & datos centrados (una fila por punto) y su covarianza \\
$(B,T,\dmodel)$ & forma de las activaciones: lote, longitud, dimensión del modelo \\
\bottomrule
\end{tabular}
\end{center}

Los vectores son columnas: $v=(v_1,\dots,v_n)$ denota un vector de $\R^{n}$ escrito en horizontal por comodidad, y $u\T v$ es su producto escalar.


\tableofcontents

\mainmatter
% =====================================================================
\chapter{Independencia lineal, bases y rango}
\label{tema:1}
% =====================================================================

Este tema introduce los conceptos que permiten medir cuántas dimensiones utiliza realmente una matriz: independencia lineal, base, dimensión, rango y núcleo. La presentación sigue los §3.1, §3.4 y~§3.5 de Strang~\cite{strang-ila}, las clases~9 a~11 del curso 18.06~\cite{mit1806} y los §2.5 a~§2.7 de~\cite{mml}. Se trabaja siempre en $\R^{n}$, aunque las definiciones valen en cualquier espacio vectorial real.

\section{Combinaciones lineales y subespacios}

\begin{definicion}[Subespacio]
Un subconjunto $W\subseteq\R^{n}$ es un \emph{subespacio vectorial} si contiene el vector nulo y es cerrado para la suma y el producto por escalares: si $u,v\in W$ y $c\in\R$, entonces $u+v\in W$ y $cu\in W$.
\end{definicion}

\begin{definicion}[Subespacio generado]
Dados $v_1,\dots,v_k\in\R^{n}$, una \emph{combinación lineal} de ellos es un vector $c_1v_1+\dots+c_kv_k$ con $c_i\in\R$. El conjunto de todas las combinaciones lineales es un subespacio, que se denota $\gen\{v_1,\dots,v_k\}$ y se llama \emph{subespacio generado} por los $v_i$.
\end{definicion}

\begin{definicion}[Independencia lineal]\label{def:indep}
Los vectores $v_1,\dots,v_k$ son \emph{linealmente independientes} si la única combinación lineal nula es la trivial:
\[
c_1v_1+\dots+c_kv_k=0\ \Longrightarrow\ c_1=\dots=c_k=0.
\]
En caso contrario se dicen \emph{linealmente dependientes}.
\end{definicion}

Si $V=[v_1\mid\cdots\mid v_k]\in\R^{n\times k}$ es la matriz cuyas columnas son los vectores, la condición equivale a que el sistema homogéneo $Vc=0$ tenga solo la solución $c=0$.

\begin{proposicion}\label{prop:dependencia}
Los vectores $v_1,\dots,v_k$ ($k\ge 2$) son linealmente dependientes si y solo si alguno de ellos es combinación lineal de los demás.
\end{proposicion}

\section{Bases, dimensión y coordenadas}

\begin{definicion}[Base]
Una \emph{base} de un subespacio $W$ es una familia de vectores de $W$ linealmente independiente que genera $W$.
\end{definicion}

\begin{proposicion}[Coordenadas]\label{prop:coordenadas}
Si $\{u_1,\dots,u_k\}$ es una base de $W$, todo $w\in W$ se escribe de forma única como $w=c_1u_1+\dots+c_ku_k$. Los números $c_1,\dots,c_k$ son las \emph{coordenadas} de $w$ en esa base.
\end{proposicion}

\begin{lema}\label{lem:steinitz}
Si $W=\gen\{v_1,\dots,v_m\}$ y $w_1,\dots,w_p\in W$ con $p>m$, entonces $w_1,\dots,w_p$ son linealmente dependientes.
\end{lema}

\begin{teorema}[Dimensión]\label{teo:dimension}
Todas las bases de un subespacio $W$ tienen el mismo número de elementos. Ese número se llama \emph{dimensión} de $W$ y se denota $\dim W$. En particular, $\dim\R^{n}=n$.
\end{teorema}

\addtocounter{definicion}{1}%

\begin{observacion}
Las coordenadas de un embedding dependen de la base en la que se expresan. Un cambio de base invertible del espacio de activaciones de un modelo, compensado en las capas siguientes, no altera sus predicciones; por eso las coordenadas individuales de un embedding no tienen, en general, un significado propio. Esta observación es el punto de partida de los métodos de interpretabilidad del módulo~23, que buscan direcciones del espacio y no coordenadas.
\end{observacion}

\section{Espacio columna, núcleo y rango}

\begin{definicion}[Espacio columna y núcleo]
Sea $A\in\R^{m\times n}$. El \emph{espacio columna} (o imagen) de $A$ es $\Col(A)=\{Ax : x\in\R^{n}\}\subseteq\R^{m}$, el subespacio generado por sus columnas. El \emph{núcleo} de $A$ es $\Nuc(A)=\{x\in\R^{n}: Ax=0\}\subseteq\R^{n}$.
\end{definicion}

\begin{definicion}[Rango]
El \emph{rango} de $A$ es $\rango(A)=\dim\Col(A)$, el número máximo de columnas linealmente independientes de~$A$.
\end{definicion}

\begin{teorema}[Rango de filas igual a rango de columnas]\label{teo:filas-columnas}
Para toda $A\in\R^{m\times n}$, $\rango(A)=\rango(A\T)$: el número máximo de filas independientes coincide con el de columnas independientes.
\end{teorema}

La factorización $A=CR$ de la demostración es la que Strang usa como punto de partida en~\cite[§1.4]{strang-ila}.

\begin{teorema}[Rango y nulidad]\label{teo:rango-nulidad}
Para toda $A\in\R^{m\times n}$,
\[
\rango(A)+\dim\Nuc(A)=n.
\]
\end{teorema}

\addtocounter{definicion}{1}%

\section{Matrices de rango 1 y rango de un producto}

\begin{proposicion}\label{prop:rango1}
Si $u\in\R^{m}$ y $v\in\R^{n}$ son no nulos, la matriz $uv\T\in\R^{m\times n}$ tiene rango~1. Recíprocamente, toda matriz de rango~1 es de esta forma.
\end{proposicion}

\begin{proposicion}[Rango de un producto]\label{prop:rango-producto}
Si $B\in\R^{m\times r}$ y $A\in\R^{r\times n}$, entonces $\rango(BA)\le\min\{\rango(A),\rango(B)\}\le r$.
\end{proposicion}

La demostración se propone en el \cref{eje:1-producto}. De la factorización $A=CR$ del \cref{teo:filas-columnas} se deduce además que toda matriz de rango $r$ es suma de $r$ matrices de rango~1, $A=\sum_{i=1}^{r}c_i\,\rho_i\T$, donde $c_i$ es la columna $i$ de $C$ y $\rho_i\T$ la fila $i$ de $R$. El tema~4 obtiene una descomposición de este tipo con piezas ortogonales y ordenadas por importancia.

\begin{observacion}[Matrices de bajo rango en un LLM]\label{obs:lora-rango}
El ajuste fino LoRA~\cite{hu-lora}, que se estudia en el módulo~20, mantiene congelada una matriz de pesos $W_0\in\R^{d_{\text{out}}\times d_{\text{in}}}$ y aprende una corrección $\Delta W=BA$ con $B\in\R^{d_{\text{out}}\times r}$ y $A\in\R^{r\times d_{\text{in}}}$. Por la \cref{prop:rango-producto}, $\rango(\Delta W)\le r$, y el número de parámetros entrenables es $r(d_{\text{in}}+d_{\text{out}})$ en lugar de $d_{\text{in}}d_{\text{out}}$. Con $d_{\text{in}}=d_{\text{out}}=4096$ y $r=8$ se entrenan $65\,536$ parámetros, el $0{,}39\,\%$ de los $16\,777\,216$ de la matriz completa.
\end{observacion}


% =====================================================================
\chapter{Inversa, determinante y matrices ortogonales}
\label{tema:2}
% =====================================================================

Este tema estudia cuándo una transformación lineal puede deshacerse, cómo lo detecta el determinante y por qué las matrices ortogonales, cuya inversa es su traspuesta, ocupan un lugar privilegiado. Se sigue a Strang~\cite[§2.2, cap.~5 y §4.4]{strang-ila}, las clases~3, 17, 18 y~20 de 18.06~\cite{mit1806} y el §4.1 de~\cite{mml}.

\section{La matriz inversa}

\begin{definicion}[Inversa]
Una matriz cuadrada $A\in\R^{n\times n}$ es \emph{invertible} si existe $A^{-1}\in\R^{n\times n}$ tal que $AA^{-1}=A^{-1}A=I_n$. Una matriz no invertible se llama \emph{singular}.
\end{definicion}

\begin{proposicion}\label{prop:inversa-props}
Sean $A,B\in\R^{n\times n}$ invertibles.
\begin{enumerate}
  \item La inversa es única.
  \item $AB$ es invertible y $(AB)^{-1}=B^{-1}A^{-1}$.
  \item $A\T$ es invertible y $(A\T)^{-1}=(A^{-1})\T$.
\end{enumerate}
\end{proposicion}

\begin{teorema}[Caracterizaciones de la invertibilidad]\label{teo:invertible}
Para $A\in\R^{n\times n}$ son equivalentes:
\begin{enumerate}
  \item $A$ es invertible;
  \item $\rango(A)=n$;
  \item $\Nuc(A)=\{0\}$, es decir, $Ax=0$ solo si $x=0$;
  \item $\det(A)\ne 0$;
  \item $0$ no es autovalor de $A$.
\end{enumerate}
\end{teorema}

\begin{proposicion}[Inversa de orden 2]\label{prop:inversa2}
Sea $A=\begin{pmatrix}a&b\\c&d\end{pmatrix}$ con $ad-bc\ne 0$. Entonces
\[
A^{-1}=\frac{1}{ad-bc}\begin{pmatrix}d&-b\\-c&a\end{pmatrix}.
\]
\end{proposicion}

\addtocounter{definicion}{1}%

\begin{observacion}[Cálculo numérico]
Para resolver $Ax=b$ no se calcula $A^{-1}$: se factoriza $A$ (por ejemplo, $A=LU$) y se resuelven dos sistemas triangulares, lo que es más rápido y numéricamente más estable. En NumPy, \texttt{np.linalg.solve(A, b)} es preferible a \texttt{np.linalg.inv(A) @ b}. En los modelos de lenguaje las matrices de pesos no se invierten; lo que sí aparece constantemente es la ortogonalidad (\cref{sec:ortogonales}).
\end{observacion}

\section{El determinante}\label{sec:determinante}

Se adopta la presentación axiomática de~\cite[§5.1]{strang-ila} y de la clase~18 de~\cite{mit1806}.

\begin{definicion}[Determinante]\label{def:det}
El \emph{determinante} es la única función $\det\colon\R^{n\times n}\to\R$ que cumple:
\begin{enumerate}
  \item $\det(I_n)=1$;
  \item cambia de signo al intercambiar dos filas;
  \item es lineal en cada fila cuando las demás se mantienen fijas.
\end{enumerate}
\end{definicion}

De estas tres propiedades se deducen todas las demás~\cite[§5.1]{strang-ila}: si dos filas son iguales, el determinante es nulo; restar a una fila un múltiplo de otra no lo cambia; el determinante de una matriz triangular es el producto de su diagonal; $\det(A\T)=\det(A)$; y
\[
\det(AB)=\det(A)\det(B).
\]
La eliminación gaussiana lleva $A$ a una matriz triangular $U$ mediante operaciones que, a lo sumo, cambian el signo del determinante. Por tanto, $\det(A)\ne 0$ si y solo si todos los pivotes son no nulos, es decir, si y solo si $\rango(A)=n$. Esto completa la equivalencia (ii)$\Leftrightarrow$(iv) del \cref{teo:invertible}. Además, de $AA^{-1}=I$ se obtiene $\det(A^{-1})=1/\det(A)$.

\begin{proposicion}[Desarrollo por cofactores]\label{prop:laplace}
Para $A\in\R^{n\times n}$ y cualquier fila $i$,
\[
\det(A)=\sum_{j=1}^{n}(-1)^{i+j}a_{ij}\det(A_{ij}),
\]
donde $A_{ij}$ es la submatriz de orden $n-1$ obtenida al suprimir la fila $i$ y la columna $j$. En particular, $\det\begin{pmatrix}a&b\\c&d\end{pmatrix}=ad-bc$.
\end{proposicion}

Esta fórmula es el desarrollo de Laplace~\cite[teorema~4.2]{mml}; su demostración, a partir de la linealidad en la fila~$i$, está en~\cite[§5.2]{strang-ila}. Para $n=3$ equivale a la regla de Sarrus.

\begin{observacion}[Interpretación geométrica]
$\abs{\det(A)}$ es el factor por el que $A$ multiplica volúmenes: el cubo unidad se transforma en un paralelepípedo de volumen $\abs{\det(A)}$~\cite[§5.3]{strang-ila}. Un determinante nulo indica que el paralelepípedo es plano, es decir, que $A$ aplasta alguna dimensión. Para la matriz $M$ del \cref{ej:M}, $\det(M)=0$.
\end{observacion}

\section{Bases ortonormales y matrices ortogonales}\label{sec:ortogonales}

\begin{definicion}[Familia ortonormal]
Los vectores $q_1,\dots,q_k\in\R^{n}$ son \emph{ortonormales} si $q_i\T q_j=\delta_{ij}$, donde $\delta_{ij}$ vale~1 si $i=j$ y~0 en otro caso.
\end{definicion}

Una familia ortonormal es linealmente independiente: si $\sum_i c_iq_i=0$, multiplicando por $q_j\T$ se obtiene $c_j=0$.

\begin{proposicion}[Coordenadas en una base ortonormal]\label{prop:coord-ortonormal}
Si $\{q_1,\dots,q_n\}$ es una base ortonormal de $\R^{n}$, todo $x\in\R^{n}$ cumple
\[
x=\sum_{j=1}^{n}(q_j\T x)\,q_j\qquad\text{y}\qquad \norm{x}^{2}=\sum_{j=1}^{n}(q_j\T x)^{2}.
\]
Si la base es solo ortogonal (con $q_j\ne 0$), la coordenada $j$ es $c_j=q_j\T x/q_j\T q_j$.
\end{proposicion}

La segunda igualdad es el teorema de Pitágoras en $\R^{n}$. Calcular coordenadas en una base ortonormal no requiere resolver ningún sistema: basta un producto escalar por coordenada.

\begin{definicion}[Matriz ortogonal]
Una matriz $Q\in\R^{n\times n}$ es \emph{ortogonal} si sus columnas forman una base ortonormal de $\R^{n}$, es decir, si $Q\T Q=I_n$. En ese caso, $Q^{-1}=Q\T$.
\end{definicion}

\addtocounter{definicion}{1}%

\begin{proposicion}\label{prop:ortogonal-isometria}
Si $Q$ es ortogonal, entonces $(Qx)\T(Qy)=x\T y$ y $\norm{Qx}=\norm{x}$ para todos $x,y$, y $\det(Q)=\pm 1$.
\end{proposicion}

La demostración se propone en el \cref{eje:2-isometria}. Las matrices ortogonales son, por tanto, las transformaciones lineales que conservan longitudes y ángulos: giros y reflexiones.

\begin{observacion}[Gram--Schmidt]
A partir de cualquier base $a_1,\dots,a_n$ se obtiene una base ortonormal del mismo espacio restando a cada $a_k$ su proyección sobre los vectores ya construidos y normalizando: $w_k=a_k-\sum_{i<k}(q_i\T a_k)\,q_i$ y $q_k=w_k/\norm{w_k}$~\cite[§4.4]{strang-ila}. En forma matricial, el proceso da la factorización $A=QR$ con $R$ triangular superior.
\end{observacion}

\begin{observacion}[Ortogonalidad en un LLM]
Las matrices ortogonales aparecen en varios puntos de un Transformer. La codificación posicional rotatoria (RoPE, módulo~14) aplica a cada par de coordenadas de las consultas y las claves un giro $Q_\theta$ cuyo ángulo depende de la posición; como los giros conservan la norma, solo cambian los ángulos entre vectores. Las matrices $U$ y $V$ de la SVD (\cref{tema:4}) son ortogonales, y algunas inicializaciones de pesos usan matrices ortogonales aleatorias porque conservan la norma de las activaciones.
\end{observacion}


% =====================================================================
\chapter{Autovalores, autovectores y teorema espectral}
\label{tema:3}
% =====================================================================

Un autovector de una matriz cuadrada es una dirección que la matriz no gira, solo estira o encoge. Este tema define autovalores y autovectores, estudia la diagonalización y demuestra el resultado central para las aplicaciones: las matrices reales simétricas tienen autovalores reales y una base ortonormal de autovectores. Se sigue a Strang~\cite[§6.1--6.3]{strang-ila}, las clases~21, 22 y~25 de 18.06~\cite{mit1806}, las clases~4 y~5 de 18.065~\cite{mit18065} y los §4.2 y~§4.4 de~\cite{mml}.

\section{Definición y polinomio característico}

\begin{definicion}[Autovalor y autovector]
Sea $A\in\R^{n\times n}$. Un número $\lambda\in\mathbb{C}$ es un \emph{autovalor} de $A$ si existe $v\ne 0$ tal que
\[
Av=\lambda v.
\]
Todo $v\ne 0$ que cumple la igualdad es un \emph{autovector} asociado a $\lambda$. El conjunto $\Nuc(A-\lambda I)$, formado por los autovectores y el vector nulo, es el \emph{espacio propio} de $\lambda$.
\end{definicion}

\begin{proposicion}[Polinomio característico]\label{prop:polcar}
$\lambda$ es autovalor de $A$ si y solo si $\det(A-\lambda I)=0$. La función $p_A(\lambda)=\det(A-\lambda I)$ es un polinomio de grado $n$ en $\lambda$, llamado \emph{polinomio característico}.
\end{proposicion}

Por el teorema fundamental del álgebra, $p_A$ tiene $n$ raíces complejas contadas con multiplicidad: $p_A(\lambda)=\prod_{i=1}^{n}(\lambda_i-\lambda)$.

\addtocounter{definicion}{1}%

\begin{proposicion}[Traza y determinante]\label{prop:traza-det}
Si $\lambda_1,\dots,\lambda_n$ son los autovalores de $A$ contados con multiplicidad, entonces
\[
\det(A)=\prod_{i=1}^{n}\lambda_i,\qquad \traza(A)=\sum_{i=1}^{n}a_{ii}=\sum_{i=1}^{n}\lambda_i.
\]
\end{proposicion}

En el \cref{ej:S}, $3+1=4=\traza(S)$ y $3\cdot 1=3=\det(S)$. Estas dos comprobaciones detectan la mayoría de los errores de cálculo.

\begin{table}[h]
\centering\small
\caption{Autovalores de algunas matrices de orden 2.}\label{tab:autovalores}
\begin{tabular}{@{}lcl@{}}
\toprule
Matriz & Autovalores & Interpretación \\
\midrule
$\begin{pmatrix}2&1\\1&2\end{pmatrix}$ & $3$ y $1$ & estira por~3 la diagonal y conserva la antidiagonal \\[6pt]
$\begin{pmatrix}3&1\\0&2\end{pmatrix}$ & $3$ y $2$ & triangular: los autovalores son la diagonal \\[6pt]
$\begin{pmatrix}1&0\\0&0\end{pmatrix}$ & $1$ y $0$ & proyección: conserva el eje $x$ y anula el eje $y$ \\[6pt]
$\begin{pmatrix}0&-1\\1&0\end{pmatrix}$ & $\pm i$ & giro de $90^{\circ}$: ninguna dirección real se conserva \\
\bottomrule
\end{tabular}
\end{table}

\section{Diagonalización}

\begin{proposicion}\label{prop:indep-autovectores}
Autovectores asociados a autovalores distintos son linealmente independientes.
\end{proposicion}

\begin{teorema}[Diagonalización]\label{teo:diagonalizacion}
$A\in\R^{n\times n}$ admite una factorización $A=X\Lambda X^{-1}$, con $\Lambda$ diagonal y $X$ invertible, si y solo si $A$ tiene $n$ autovectores linealmente independientes. En ese caso, las columnas de $X$ son autovectores, la diagonal de $\Lambda$ contiene los autovalores correspondientes y $A^{k}=X\Lambda^{k}X^{-1}$ para todo $k\ge 1$.
\end{teorema}

Por la \cref{prop:indep-autovectores}, una matriz con $n$ autovalores distintos es diagonalizable. Con autovalores repetidos puede no serlo: $\begin{pmatrix}1&1\\0&1\end{pmatrix}$ tiene el autovalor doble~1, pero su espacio propio es $\gen\{e_1\}$, de dimensión~1. La diagonalizabilidad no guarda relación con la invertibilidad (\cref{eje:3-mml}).

\section{Matrices simétricas y teorema espectral}

\begin{teorema}\label{teo:simetricas}
Sea $S\in\R^{n\times n}$ simétrica, $S\T=S$. Entonces:
\begin{enumerate}
  \item todos los autovalores de $S$ son reales;
  \item autovectores asociados a autovalores distintos son ortogonales.
\end{enumerate}
\end{teorema}

\begin{teorema}[Teorema espectral]\label{teo:espectral}
Toda matriz real simétrica $S\in\R^{n\times n}$ se factoriza como
\[
S=Q\Lambda Q\T=\sum_{i=1}^{n}\lambda_i\,q_iq_i\T,
\]
donde $Q=[q_1\mid\cdots\mid q_n]$ es ortogonal y $\Lambda=\diag(\lambda_1,\dots,\lambda_n)$ es real.
\end{teorema}

Leída de derecha a izquierda, $S=Q\Lambda Q\T$ dice que $S$ expresa el vector en la base de autovectores ($Q\T$), estira cada eje por su autovalor ($\Lambda$) y vuelve a la base original ($Q$). Para el \cref{ej:S},
\[
S=3\,q_1q_1\T+1\,q_2q_2\T=\frac{3}{2}\begin{pmatrix}1&1\\1&1\end{pmatrix}+\frac{1}{2}\begin{pmatrix}1&-1\\-1&1\end{pmatrix}.
\]

\begin{definicion}[Matriz semidefinida positiva]
Una matriz simétrica $S$ es \emph{semidefinida positiva} si $x\T Sx\ge 0$ para todo $x$, y \emph{definida positiva} si $x\T Sx>0$ para todo $x\ne 0$.
\end{definicion}

\begin{proposicion}\label{prop:AtA}
Para toda $A\in\R^{m\times n}$, la matriz $A\T A$ es simétrica y semidefinida positiva, y sus autovalores son mayores o iguales que cero.
\end{proposicion}

\begin{proposicion}[Cociente de Rayleigh]\label{prop:rayleigh}
Sea $S$ simétrica con autovalores $\lambda_1\ge\dots\ge\lambda_n$ y autovectores ortonormales $q_1,\dots,q_n$. Entonces
\[
\max_{\norm{x}=1}x\T Sx=\lambda_1,
\]
y el máximo se alcanza en $x=q_1$.
\end{proposicion}

Este resultado es la base del análisis de componentes principales (\cref{tema:5}): la dirección de máxima varianza de unos datos es el primer autovector de su matriz de covarianza.

\section{El método de la potencia}

Si $S$ es simétrica con $\abs{\lambda_1}>\abs{\lambda_2}\ge\dots\ge\abs{\lambda_n}$, la iteración
\[
x_{k+1}=\frac{Sx_k}{\norm{Sx_k}}
\]
converge, salvo signo, al autovector $q_1$ siempre que $x_0$ tenga componente no nula sobre $q_1$. El error decrece como $\abs{\lambda_2/\lambda_1}^{k}$ (\cref{eje:3-potencia}).

\addtocounter{definicion}{1}%

\addtocounter{definicion}{1}%


% =====================================================================
\chapter{La descomposición en valores singulares}
\label{tema:4}
% =====================================================================

El teorema espectral solo se aplica a matrices cuadradas simétricas. La descomposición en valores singulares (SVD) extiende la misma idea a cualquier matriz rectangular: toda transformación lineal es un giro, seguido de un estiramiento a lo largo de ejes perpendiculares, seguido de otro giro. Este tema sigue a Strang~\cite[§7.1]{strang-ila}, la clase~29 de 18.06~\cite{mit1806}, las clases~6 y~8 de 18.065~\cite{mit18065} y el §4.5 de~\cite{mml}.

\section{El teorema de la SVD}

\begin{teorema}[Descomposición en valores singulares]\label{teo:svd}
Sea $A\in\R^{m\times n}$ de rango $r$. Existen matrices ortogonales $U\in\R^{m\times m}$ y $V\in\R^{n\times n}$ y números $\sigma_1\ge\sigma_2\ge\dots\ge\sigma_r>0$ tales que
\[
A=U\Sigma V\T=\sum_{i=1}^{r}\sigma_i\,u_iv_i\T,
\]
donde $\Sigma\in\R^{m\times n}$ tiene $\Sigma_{ii}=\sigma_i$ para $i\le r$ y ceros en el resto, y $u_i$, $v_i$ son las columnas de $U$ y $V$. Los $\sigma_i$ son los \emph{valores singulares} de $A$, y los $u_i$ y $v_i$, sus \emph{vectores singulares} izquierdos y derechos. Además, $Av_i=\sigma_iu_i$ y $A\T u_i=\sigma_iv_i$ para $i\le r$.
\end{teorema}

\addtocounter{figure}{1}%

\begin{observacion}[SVD reducida]
Como las columnas $u_i$ y $v_i$ con $i>r$ se multiplican por ceros, también se cumple $A=U_r\Sigma_rV_r\T$, con $U_r\in\R^{m\times r}$, $\Sigma_r=\diag(\sigma_1,\dots,\sigma_r)$ y $V_r\in\R^{n\times r}$. Es la versión que se almacena en la práctica (\texttt{np.linalg.svd} con \texttt{full\_matrices=False}).
\end{observacion}

\begin{corolario}\label{cor:svd-props}
Sea $A=U\Sigma V\T$ la SVD de $A\in\R^{m\times n}$.
\begin{enumerate}
  \item $\rango(A)$ es el número de valores singulares no nulos.
  \item $\sigma_i(A)=\sqrt{\lambda_i(A\T A)}$, y $AA\T$ tiene los mismos autovalores no nulos que $A\T A$, con autovectores $u_i$.
  \item $\{u_1,\dots,u_r\}$ es una base ortonormal de $\Col(A)$ y $\{v_{r+1},\dots,v_n\}$, una base ortonormal de $\Nuc(A)$.
  \item Si $A$ es simétrica con autovalores $\lambda_i$, sus valores singulares son los $\abs{\lambda_i}$. Si además es semidefinida positiva, la SVD coincide con la descomposición espectral.
\end{enumerate}
\end{corolario}

\section{Cálculo a mano}

El procedimiento de la demostración sirve también para calcular la SVD de matrices pequeñas: (1) se forma $A\T A$ o $AA\T$, la que sea de menor orden; (2) se calculan sus autovalores $\sigma_i^{2}$ y autovectores ortonormales; (3) se obtienen los vectores singulares del otro lado con $u_i=Av_i/\sigma_i$ o $v_i=A\T u_i/\sigma_i$.

\addtocounter{definicion}{1}%

\section{Normas matriciales}

\begin{definicion}[Normas de Frobenius y espectral]
Para $A\in\R^{m\times n}$ se definen la \emph{norma de Frobenius} y la \emph{norma espectral}
\[
\norm{A}_F=\Bigl(\sum_{i=1}^{m}\sum_{j=1}^{n}a_{ij}^{2}\Bigr)^{1/2},\qquad
\norm{A}_2=\max_{x\ne 0}\frac{\norm{Ax}}{\norm{x}}.
\]
\end{definicion}

\begin{proposicion}\label{prop:normas-svd}
Para toda $A\in\R^{m\times n}$ con valores singulares $\sigma_1\ge\dots\ge\sigma_r>0$,
\[
\norm{A}_F^{2}=\traza(A\T A)=\sum_{i=1}^{r}\sigma_i^{2},\qquad \norm{A}_2=\sigma_1.
\]
En particular, $\norm{A}_2\le\norm{A}_F\le\sqrt{r}\,\norm{A}_2$.
\end{proposicion}

Las dos igualdades se demuestran en los \cref{eje:4-frobenius,eje:4-espectral}; la segunda es el teorema~4.24 de~\cite{mml}. Ambas normas son invariantes por transformaciones ortogonales: $\norm{Q_1AQ_2}=\norm{A}$ para $Q_1$, $Q_2$ ortogonales, porque los valores singulares no cambian. Las desigualdades finales se siguen de $\sigma_1^{2}\le\sum_i\sigma_i^{2}\le r\sigma_1^{2}$.

\addtocounter{definicion}{1}%

\addtocounter{definicion}{1}%


% =====================================================================
\chapter{Aproximación de bajo rango y análisis de componentes principales}
\label{tema:5}
% =====================================================================

La SVD escribe una matriz como suma de piezas de rango~1 ordenadas por su valor singular. Este tema demuestra que quedarse con las $k$ primeras piezas da la mejor aproximación posible de rango $k$ (teorema de Eckart--Young~\cite{eckart-young}) y extrae tres consecuencias: la compresión de matrices, el ajuste fino de bajo rango LoRA~\cite{hu-lora} y el análisis de componentes principales (PCA). Se sigue la clase~7 de 18.065~\cite{mit18065}, los §7.2 y~§7.3 de~\cite{strang-ila}, el §4.6 y el cap.~10 de~\cite{mml} y el tutorial de Shlens~\cite{shlens}.

\section{El teorema de Eckart--Young}

Sea $A=\sum_{i=1}^{r}\sigma_iu_iv_i\T$ la SVD de $A\in\R^{m\times n}$ y, para $k<r$, sea
\[
A_k=\sum_{i=1}^{k}\sigma_i\,u_iv_i\T,
\]
que tiene rango $k$ y cuya SVD es la de $A$ truncada. Como $A-A_k=\sum_{i>k}\sigma_iu_iv_i\T$ es a su vez una SVD, la \cref{prop:normas-svd} da
\[
\norm{A-A_k}_2=\sigma_{k+1},\qquad \norm{A-A_k}_F=\Bigl(\sum_{i=k+1}^{r}\sigma_i^{2}\Bigr)^{1/2}.
\]

\begin{lema}\label{lem:interseccion}
Si $W_1$ y $W_2$ son subespacios de $\R^{n}$ con $\dim W_1+\dim W_2>n$, entonces $W_1\cap W_2$ contiene un vector no nulo.
\end{lema}

\begin{teorema}[Eckart--Young]\label{teo:eckart-young}
Sea $A\in\R^{m\times n}$ con valores singulares $\sigma_1\ge\sigma_2\ge\cdots$ y sea $k<\rango(A)$. Para toda $B\in\R^{m\times n}$ con $\rango(B)\le k$,
\[
\norm{A-B}_2\ge\sigma_{k+1}=\norm{A-A_k}_2
\qquad\text{y}\qquad
\norm{A-B}_F\ge\Bigl(\sum_{i>k}\sigma_i^{2}\Bigr)^{1/2}=\norm{A-A_k}_F.
\]
Es decir, $A_k$ es la mejor aproximación de rango a lo sumo $k$ de $A$ en ambas normas.
\end{teorema}

La parte espectral es el teorema~4.25 de~\cite{mml}; la de Frobenius es el enunciado original de Eckart y Young~\cite{eckart-young}. La desigualdad $\sigma_{k+j}(A)\le\sigma_j(A-B)$ usada en la segunda parte es un caso de la desigualdad de Weyl para valores singulares.

\addtocounter{definicion}{1}%

\section{Compresión de matrices}

Almacenar $A_k$ requiere $k$ columnas de $U$, $k$ filas de $V\T$ y $k$ valores singulares: $k(m+n+1)$ números frente a los $mn$ de $A$. La compresión compensa cuando $k(m+n+1)<mn$, y es eficaz cuando los valores singulares decrecen deprisa, pues entonces el error $(\sum_{i>k}\sigma_i^{2})^{1/2}$ es pequeño con $k$ pequeño.

\begin{table}[h]
\centering\small
\caption{Coste de almacenamiento de la aproximación de rango $k$.}\label{tab:compresion}
\begin{tabular}{@{}lrrrr@{}}
\toprule
Matriz & $mn$ & $k$ & $k(m+n+1)$ & Fracción \\
\midrule
Imagen de $512\times 512$ & $262\,144$ & 20 & $20\,500$ & $7{,}8\,\%$ \\
Imagen de $1000\times 800$ & $800\,000$ & 50 & $90\,050$ & $11{,}3\,\%$ \\
Pesos de $4096\times 4096$ & $16\,777\,216$ & 8 & $65\,544$ & $0{,}39\,\%$ \\
\bottomrule
\end{tabular}
\end{table}

Una imagen en escala de grises es una matriz de intensidades. Las regiones uniformes, los bordes horizontales y verticales y los degradados suaves tienen rango bajo; los bordes diagonales y el ruido requieren muchos términos~\cite[§7.2]{strang-ila}. La práctica del módulo aplica la SVD a una imagen sintética y muestra cómo decae el error con $k$.

\section{Ajuste fino de bajo rango}

\addtocounter{figure}{1}%

LoRA~\cite{hu-lora} parte de la hipótesis de que la modificación que el ajuste fino introduce en los pesos de un modelo preentrenado tiene un \emph{rango intrínseco} bajo. En lugar de actualizar $W_0$, aprende $\Delta W=BA$ (\cref{fig:lora}), con $r(d_{\text{in}}+d_{\text{out}})$ parámetros (\cref{obs:lora-rango}). Al terminar, $BA$ puede sumarse a $W_0$, de modo que la inferencia no tiene coste adicional. El \cref{teo:eckart-young} cuantifica lo que se pierde: si la actualización ideal es $\Delta W^{\star}$, ninguna corrección de rango $r$ puede acercarse a ella más que $(\sum_{i>r}\sigma_i^{2}(\Delta W^{\star}))^{1/2}$ (\cref{eje:5-lora}). La apuesta de LoRA, respaldada empíricamente en~\cite{hu-lora}, es que el espectro de $\Delta W^{\star}$ decae deprisa.

\section{Análisis de componentes principales}

\begin{definicion}[Datos centrados y covarianza]
Sean $n$ puntos $x_1,\dots,x_n\in\R^{d}$ con media $\mu=\frac1n\sum_ix_i$. La matriz de datos centrados $X\in\R^{n\times d}$ tiene por filas $(x_i-\mu)\T$. La \emph{matriz de covarianza} es
\[
C=\frac{1}{n}X\T X\in\R^{d\times d}.
\]
\end{definicion}

$C$ es simétrica y semidefinida positiva (\cref{prop:AtA}). Algunos textos dividen por $n-1$ en lugar de $n$; el factor no cambia los autovectores. Para un vector unitario $w$, las proyecciones de los puntos sobre $w$ son las componentes de $Xw$, que tienen media cero, y su varianza es
\[
\frac1n\norm{Xw}^{2}=\frac1n w\T X\T Xw=w\T Cw.
\]

\begin{teorema}[Dirección de máxima varianza]\label{teo:pca}
Sean $\lambda_1\ge\dots\ge\lambda_d\ge 0$ los autovalores de $C$ y $q_1,\dots,q_d$ autovectores ortonormales. La dirección unitaria $w$ que maximiza la varianza de las proyecciones es $w=q_1$, y la varianza máxima es $\lambda_1$. Más en general, entre todos los subespacios de dimensión $k$, el que maximiza la varianza total de las proyecciones es $\gen\{q_1,\dots,q_k\}$, con varianza $\lambda_1+\dots+\lambda_k$.
\end{teorema}

Los vectores $q_i$ son las \emph{componentes principales}, y la \emph{varianza explicada} por la componente $i$ es $\lambda_i/\sum_j\lambda_j$.

\begin{teorema}[PCA mediante la SVD]\label{teo:svd-pca}
Sea $X=U\Sigma V\T$ la SVD de la matriz de datos centrados. Entonces las componentes principales son las columnas de $V$, los autovalores de la covarianza son $\lambda_i=\sigma_i^{2}/n$ y las coordenadas de los datos en las $k$ primeras componentes son
\[
Z_k=XV_k=U_k\Sigma_k\in\R^{n\times k},
\]
cuya columna $i$ tiene varianza $\lambda_i$.
\end{teorema}

La reconstrucción de los datos a partir de $k$ componentes es $Z_kV_k\T=U_k\Sigma_kV_k\T=X_k$, el truncamiento de la SVD de $X$. Por el \cref{teo:eckart-young}, PCA es la mejor aproximación de rango $k$ de los datos centrados, con error cuadrático total $\sum_{i>k}\sigma_i^{2}$: maximizar la varianza proyectada y minimizar el error de reconstrucción son el mismo problema~\cite[§10.3]{mml}.

\addtocounter{definicion}{1}%

\begin{observacion}[Cálculo numérico]
En la práctica se calcula la SVD de $X$ y no la descomposición espectral de $X\T X$: formar $X\T X$ eleva al cuadrado el número de condición $\sigma_1/\sigma_r$ y pierde precisión en las componentes de varianza pequeña~\cite[§10.4]{mml}.
\end{observacion}

\begin{observacion}[Visualización de embeddings]
Los embeddings de GPT-2 small tienen $\dmodel=768$ componentes~\cite{gpt2}. Proyectarlos sobre las dos primeras componentes principales da la imagen plana que conserva más varianza, y palabras de un mismo campo semántico suelen aparecer agrupadas. Conviene leer siempre la varianza explicada: una proyección ortogonal nunca alarga distancias, de modo que dos puntos alejados en el plano también lo están en $\R^{768}$, pero dos puntos cercanos en el plano pueden estar lejos en el espacio original (\cref{eje:5-varexp}). Los métodos de interpretabilidad del módulo~23 usan PCA para encontrar direcciones dominantes en las activaciones de un modelo.
\end{observacion}


% =====================================================================
\chapter{Tensores, \emph{broadcasting} y formas de un LLM}
\label{tema:6}
% =====================================================================

En el código de un modelo de lenguaje casi nunca aparecen vectores o matrices aislados, sino arrays de varios ejes que agrupan muchos vectores a la vez. Este tema formaliza esos arrays, el producto matricial por lotes y las reglas de \emph{broadcasting}, y las aplica a las formas de las matrices de un Transformer y al rango de su capa de salida. Se sigue a Zhang \emph{et al.}~\cite[§2.1 y §2.3]{d2l}, la documentación de NumPy~\cite{numpy-bc} y el apéndice~6 de~\cite{strang-ila}; la última sección, a Yang \emph{et al.}~\cite{yang-softmax}.

\section{Tensores como arrays}

\begin{definicion}[Tensor]
Un \emph{tensor} de orden $k$ y forma $(n_1,\dots,n_k)$ es una familia de números reales $X[i_1,\dots,i_k]$ indexada por $0\le i_j<n_j$; contiene $n_1n_2\cdots n_k$ elementos. Cada índice recorre un \emph{eje}. Los escalares, vectores y matrices son tensores de orden 0, 1 y~2.
\end{definicion}

En física y en geometría diferencial, un tensor es una aplicación multilineal cuyas componentes cambian de una forma determinada al cambiar de base~\cite[apéndice~6]{strang-ila}. En aprendizaje automático el término designa simplemente un array multidimensional~\cite[§2.3.4]{d2l}, y es el sentido que se usa en este curso.

\begin{definicion}[Reducción a lo largo de un eje]
Sumar un tensor de forma $(n_1,\dots,n_k)$ a lo largo del eje $j$ produce un tensor de forma $(n_1,\dots,n_{j-1},n_{j+1},\dots,n_k)$:
\[
Y[i_1,\dots,\widehat{i_j},\dots,i_k]=\sum_{i_j=0}^{n_j-1}X[i_1,\dots,i_k],
\]
donde $\widehat{i_j}$ indica que el índice se suprime. Con la opción \texttt{keepdims=True}, el eje $j$ se conserva con tamaño~1.
\end{definicion}

La media, el máximo y la función $\LSE$ del módulo~1 se reducen del mismo modo. En un LLM, la softmax se aplica a lo largo del último eje de los logits y la normalización por capas calcula media y varianza a lo largo del eje de $\dmodel$.

\section{Producto por lotes y notación de Einstein}

\begin{definicion}[Producto matricial por lotes]
Sean $X$ de forma $(\dots,m,p)$ e $Y$ de forma $(\dots,p,q)$, con los mismos ejes iniciales (o ejes compatibles por \emph{broadcasting}). El producto $X\,@\,Y$ tiene forma $(\dots,m,q)$ y se obtiene multiplicando como matrices los dos últimos ejes para cada valor de los índices iniciales, que actúan como índices de lote.
\end{definicion}

Multiplicar una matriz de $m\times p$ por otra de $p\times q$ requiere $mpq$ multiplicaciones. Por eso el orden en que se agrupa un producto de varias matrices afecta mucho al coste (\cref{eje:6-orden}).

La notación de Einstein (\texttt{np.einsum}) describe estas operaciones nombrando los índices: los que aparecen en las entradas y no en la salida se suman. Por ejemplo, para $Q$ y $K$ de forma $(B,T,d)$, la expresión \texttt{einsum('bid,bjd->bij', Q, K)} calcula
\[
S[b,i,j]=\sum_{k}Q[b,i,k]\,K[b,j,k],
\]
es decir, $S=QK\T$ para cada elemento del lote, de forma $(B,T,T)$. Es el núcleo de la atención del módulo~13. Con \texttt{@}, el mismo cálculo exige trasponer solo los dos últimos ejes: \texttt{Q @ K.transpose(0, 2, 1)}.

\section{Broadcasting}

\begin{definicion}[Reglas de \emph{broadcasting}]\label{def:broadcasting}
Dos tensores de formas $(a_1,\dots,a_p)$ y $(b_1,\dots,b_q)$ se combinan elemento a elemento del siguiente modo~\cite{numpy-bc}:
\begin{enumerate}
  \item se alinean las formas por la derecha y se completa la más corta con ejes de tamaño~1 por la izquierda;
  \item en cada eje, los tamaños son \emph{compatibles} si son iguales o si uno de ellos vale~1; en otro caso la operación produce un error;
  \item el resultado tiene, en cada eje, el mayor de los dos tamaños, y el tensor con tamaño~1 en ese eje se repite a lo largo de él.
\end{enumerate}
\end{definicion}

\begin{table}[h]
\centering\small
\caption{Ejemplos de \emph{broadcasting} con activaciones de forma $(2,3,4)$.}\label{tab:broadcasting}
\begin{tabular}{@{}lll@{}}
\toprule
Operación & Formas & Resultado \\
\midrule
\texttt{x + b} & $(2,3,4)+(4,)$ & $(2,3,4)$: sesgo sumado a cada token \\
\texttt{x - x.mean(-1, keepdims=True)} & $(2,3,4)-(2,3,1)$ & $(2,3,4)$: centra cada token \\
\texttt{x @ W} & $(2,3,4)\,@\,(4,8)$ & $(2,3,8)$: misma capa en cada token \\
\texttt{x + y} & $(2,3,4)+(3,)$ & error: $4\ne 3$ \\
\bottomrule
\end{tabular}
\end{table}

\begin{observacion}
El \emph{broadcasting} es la forma eficiente de escribir operaciones que en notación matemática llevan un vector de unos: si $X\in\R^{T\times d}$ y $b\in\R^{d}$, \texttt{X + b} calcula $X+\mathbf{1}b\T$ sin construir la matriz $\mathbf{1}b\T$. Su contrapartida es que algunas combinaciones de formas no producen error y dan un resultado no deseado; el \cref{eje:6-formas} muestra un caso típico. Comprobar las formas en cada paso es el hábito que más errores evita al programar un Transformer.
\end{observacion}

\section{Formas de un LLM}

El \cref{tab:formas} recoge las formas de los tensores principales de un Transformer de tipo decodificador, con los valores de GPT-2 small~\cite{gpt2}: vocabulario de $V=50\,257$ tokens, $\dmodel=768$, 12 cabezas de atención y contexto de 1024 tokens. Se escribe $B$ para el tamaño del lote, $T$ para la longitud de la secuencia y $h$ para el número de cabezas. Las capas lineales se escriben aquí como en el código, $y=xW$ con $W$ de forma $(d_{\text{in}},d_{\text{out}})$.

\begin{table}[h]
\centering\small
\caption{Formas de los tensores de un Transformer decodificador.}\label{tab:formas}
\begin{tabular}{@{}lll@{}}
\toprule
Objeto & Forma & GPT-2 small, $B=8$ \\
\midrule
Identificadores de tokens & $(B,T)$ & $(8,1024)$ \\
Matriz de embeddings $E$ & $(V,\dmodel)$ & $(50\,257,768)$ \\
Activaciones & $(B,T,\dmodel)$ & $(8,1024,768)$ \\
Pesos de la primera capa del MLP & $(\dmodel,4\dmodel)$ & $(768,3072)$ \\
Puntuaciones de atención & $(B,h,T,T)$ & $(8,12,1024,1024)$ \\
Logits & $(B,T,V)$ & $(8,1024,50\,257)$ \\
\bottomrule
\end{tabular}
\end{table}

\addtocounter{figure}{1}%

\subsection*{La matriz de embeddings y el cuello de botella de la softmax}

La matriz de embeddings $E\in\R^{V\times\dmodel}$ tiene muchas más filas que columnas. En GPT-2 small, $50\,257\cdot 768=38\,597\,376$ parámetros, cerca del 31\,\% de los $124$~millones del modelo; en Llama~2 7B~\cite{llama2}, $32\,000\cdot 4096=131\,072\,000$. GPT-2 \emph{ata} los pesos de entrada y salida: los logits del token siguiente se calculan como $z=hE\T$, donde $h\in\R^{\dmodel}$ es la última activación.

\begin{proposicion}[Cuello de botella de la softmax]\label{prop:cuello}
Sean $N$ contextos con activaciones finales $h_1,\dots,h_N\in\R^{d}$, apiladas en $H\in\R^{N\times d}$, y sea $W\in\R^{V\times d}$ la matriz de salida. La matriz de logits $Z=HW\T\in\R^{N\times V}$ tiene rango a lo sumo $d$, y la matriz de log-probabilidades $L$, con $L_{cj}=\ln\softmax(z_c)_j$, tiene rango a lo sumo $d+1$.
\end{proposicion}

Yang \emph{et al.}~\cite{yang-softmax} plantean el modelado del lenguaje como una factorización de la matriz de log-probabilidades verdaderas. Si esa matriz, sobre todos los contextos posibles, tiene rango muy superior a $d+1$, ningún modelo con salida softmax de dimensión $d$ puede reproducirla exactamente. Con $V=50\,257$ y $d=768$, la restricción es severa en principio; los autores proponen una mezcla de varias softmax para superarla.



\backmatter
% =====================================================================
\begin{thebibliography}{99}
\addcontentsline{toc}{chapter}{Bibliografía}
\markboth{Bibliografía}{}
\small

\bibitem{mml}
M.~P. Deisenroth, A.~A. Faisal y C.~S. Ong.
\newblock \emph{Mathematics for Machine Learning}.
\newblock Cambridge University Press, 2020. Versión en línea del 15 de enero de 2024.
\newblock Cap.~2 «Linear Algebra», cap.~4 «Matrix Decompositions» y cap.~10 «Dimensionality Reduction with Principal Component Analysis».
\newblock \url{https://mml-book.github.io/}

\bibitem{eckart-young}
C.~Eckart y G.~Young.
\newblock The approximation of one matrix by another of lower rank.
\newblock \emph{Psychometrika}, 1(3):211--218, 1936.

\bibitem{gbc}
I.~Goodfellow, Y.~Bengio y A.~Courville.
\newblock \emph{Deep Learning}.
\newblock MIT Press, 2016. Cap.~2 «Linear Algebra».
\newblock \url{https://www.deeplearningbook.org/}

\bibitem{hu-lora}
E.~J. Hu, Y.~Shen, P.~Wallis, Z.~Allen-Zhu, Y.~Li, S.~Wang, L.~Wang y W.~Chen.
\newblock LoRA: Low-rank adaptation of large language models.
\newblock En \emph{International Conference on Learning Representations (ICLR)}, 2022.
\newblock \url{https://arxiv.org/abs/2106.09685}

\bibitem{numpy-bc}
NumPy Developers.
\newblock Broadcasting. \emph{NumPy User Guide}, versión~2.5.
\newblock \url{https://numpy.org/doc/stable/user/basics.broadcasting.html}

\bibitem{gpt2}
A.~Radford, J.~Wu, R.~Child, D.~Luan, D.~Amodei e I.~Sutskever.
\newblock Language models are unsupervised multitask learners.
\newblock Informe técnico, OpenAI, 2019.

\bibitem{shlens}
J.~Shlens.
\newblock A tutorial on principal component analysis.
\newblock arXiv:1404.1100, 2014.
\newblock \url{https://arxiv.org/abs/1404.1100}

\bibitem{strang-ila}
G.~Strang.
\newblock \emph{Introduction to Linear Algebra}, 6.ª ed.
\newblock Wellesley-Cambridge Press, 2023.
\newblock \url{https://math.mit.edu/~gs/linearalgebra/ila6/}

\bibitem{strang-lafd}
G.~Strang.
\newblock \emph{Linear Algebra and Learning from Data}.
\newblock Wellesley-Cambridge Press, 2019.
\newblock Texto del curso 18.065 del MIT.

\bibitem{mit1806}
G.~Strang.
\newblock 18.06: Linear Algebra.
\newblock Massachusetts Institute of Technology, MIT OpenCourseWare, primavera de 2010.
\newblock \url{https://ocw.mit.edu/courses/18-06-linear-algebra-spring-2010/}

\bibitem{mit1806sc}
G.~Strang.
\newblock 18.06SC: Linear Algebra.
\newblock Massachusetts Institute of Technology, MIT OpenCourseWare, otoño de 2011.
\newblock \url{https://ocw.mit.edu/courses/18-06sc-linear-algebra-fall-2011/}

\bibitem{mit18065}
G.~Strang.
\newblock 18.065: Matrix Methods in Data Analysis, Signal Processing, and Machine Learning.
\newblock Massachusetts Institute of Technology, MIT OpenCourseWare, primavera de 2018.
\newblock \url{https://ocw.mit.edu/courses/18-065-matrix-methods-in-data-analysis-signal-processing-and-machine-learning-spring-2018/}

\bibitem{llama2}
H.~Touvron \emph{et al.}
\newblock Llama 2: Open foundation and fine-tuned chat models.
\newblock arXiv:2307.09288, 2023.
\newblock \url{https://arxiv.org/abs/2307.09288}

\bibitem{yang-softmax}
Z.~Yang, Z.~Dai, R.~Salakhutdinov y W.~W. Cohen.
\newblock Breaking the softmax bottleneck: A high-rank RNN language model.
\newblock En \emph{International Conference on Learning Representations (ICLR)}, 2018.
\newblock \url{https://arxiv.org/abs/1711.03953}

\bibitem{d2l}
A.~Zhang, Z.~C. Lipton, M.~Li y A.~J. Smola.
\newblock \emph{Dive into Deep Learning}, versión 1.0.3.
\newblock Cambridge University Press, 2023.
\newblock \url{https://d2l.ai/}

\end{thebibliography}


\end{document}
